% ============================================================================
\clearpage
\section*{v4.9.0-r51 Source-Only: The Axis \texorpdfstring{\(e_7\)}{e7} and the Frames}
\addcontentsline{toc}{section}{v4.9.0-r51 Source-Only: The Axis e7 and the Frames}
\label{sec:v490-r51-e7-frames}
% ============================================================================

\paragraph{Purpose.}
Layer v4.9.0-r49 left open whether \(e_7\) serves at least as the reference axis of the frames.  This layer records the pre-registered test
H35 (note \texttt{llm/Vorab\_Notiz\_H35\_e7\_als\_Frame\_Achse.md}, commit \texttt{42e4fcc}, fixed before computing).  The note discloses that the
outcome was seen analytically beforehand.  The layer also records a correction to the frame rows of earlier layers; those layers are not
rewritten.

\noindent\textbf{Source-only status.}
This v4.9.0-r51 layer is source-only.  It does not build a PDF, does not assign a new DOI,
and does not authorize a public release.  The v4.9.0 Version DOI
\href{https://doi.org/10.5281/zenodo.23534954}{10.5281/zenodo.23534954} and the Concept DOI
\href{https://doi.org/10.5281/zenodo.19210728}{10.5281/zenodo.19210728} are unchanged by it.

\subsection*{v4.9.0-r51 setting}
\label{subsec:v490-r51-setting}

Under the assumption of layer v4.9.0-r33 that the \(\mathbb C\) factor of the Dixon store is the amplitude plane \(z\), the three
involutions of the doublet act on \(\mathbb C\otimes\mathbb O\) through that factor: the tachyon map of Convention~F, \(z\mapsto iz\)
(corrected in v4.9.0-r28; \(\sigma_y\) on the doublet), is multiplication \(R\) by the outer unit; the time reversal \(T\colon z\mapsto\bar z\)
(\(\sigma_z\), exchanging compensator and partner port) is the conjugation of that factor; and \(J_B=R\circ T\) (\(\sigma_x\)).  The inner unit is
\(J=L_{e_7}\) as in v4.9.0-r35.

\subsection*{v4.9.0-r51 results}
\label{subsec:v490-r51-results}

\paragraph{Theoremlet v4.9.0-r51.1.}
\begin{enumerate}[leftmargin=2em]
\item[(a)] \(R\) commutes with \(iL_{e_7}\), \(iR_{e_7}\), \(Q\), \(B-L\), the colour algebra and every sector projector \(P_s\) of v4.9.0-r48;
\(T\) and \(J_B\) map \(P_{(\ell,r)}\) to \(P_{(-\ell,-r)}\).
\item[(b)] \(L_{e_7}\) commutes with \(\operatorname{Stab}_{G_2}(e_7)\), so the tachyon map has the colour-compatible inner counterpart
\(z\mapsto e_7z\) on all of \(\mathbb O\).  The real-linear maps \(K\) of \(\mathbb O\) that commute with \(\operatorname{Stab}_{G_2}(e_7)\) and
satisfy \(KL_{e_7}=-L_{e_7}K\) form a two-dimensional space, and each vanishes on \(e_7^\perp\): \(T\) and \(J_B\) have no colour-compatible
inner counterpart on the triplet.
\item[(c)] \(\varphi(x+iy)=x+y\,e_7\) maps the frame plane onto \(\operatorname{span}(1,e_7)\) and carries \(z\mapsto iz\), \(z\mapsto\bar z\) and
\(z\mapsto i\bar z\) to \(L_{e_7}\), the octonion conjugation and their product.
\end{enumerate}

\paragraph{Proof sketch.}
(a) by matrices (script \texttt{llm/k91\_test\_h35.py}); the statements on \(R\) and \(P_s\) are also checked in Lean (\texttt{R\_iL},
\texttt{R\_iR}, \texttt{R\_P4} in \texttt{formal/QSFormal/Sektoren.lean}), those on \(T\) in v4.9.0-r48.  (b) by solving the linear conditions
with the derivation algebra of \(\mathbb O\) (dimension \(14\), stabilizer of \(e_7\) dimension \(8\)); the triplet \(3\) is not equivalent to
\(\bar 3\).  (c) since \(\operatorname{span}(1,e_7)\cong\mathbb C\).  \(\square\)

\paragraph{Verdict and reading.}
H35a, b and c are confirmed.
\begin{enumerate}[leftmargin=2em]
\item \(e_7\) is a reference axis of the frames only in part.  The tachyon map has an inner, colour-compatible counterpart \(L_{e_7}\) on all of
\(\mathbb O\); on the singlet plane all three frame maps have one, so there the frame plane is the plane of the photon kernel.  Time reversal
and \(J_B\) need the outer unit on the colour triplet.
\item Charges are frame-free between matter and tachyon: the tachyon map leaves every sector fixed, so \(Q\), \(B-L\) and the chirality are the
same in M and T.  Charge conjugation \(Q\mapsto1-Q\) is the time-reversal part \(T\), the exchange of compensator and partner port, not the
tachyon map.
\item \emph{Correction of frame rows.}  Where layers v4.9.0-r41, r43, r47, r48 and r49 call ``the frame \(T\)'' the charge conjugation, the time
reversal \(T\) is meant, not the tachyon frame of the M/T/P column; in that column the results concerned are the same in M and T.  The Lean names
\texttt{frameT} and \texttt{frameT\_swaps\_left} likewise denote the time reversal.  This agrees with v4.9.0-r34.1(d), where \(R\) fixes \(Q\)
and the effect comes from the time-reversal part.
\item The compensator \(\tfrac12(1+iL_{e_7})\) of v4.9.0-r42 is the compensator port and its complement the partner port; \(T\) exchanges the
ports and hence the chiralities, while \(R\) fixes both.  Since \(iL_{e_7}\) is the outer \(z\mapsto iz\) composed with the inner
\(z\mapsto e_7z\), the left-handed states are those on which the two tachyon maps agree up to sign.
\end{enumerate}

\subsection*{v4.9.0-r51 status in the three frames}
\label{subsec:v490-r51-frames}

\begin{center}
\small
\begin{tabular}{@{}>{\raggedright\arraybackslash}p{0.30\textwidth}>{\raggedright\arraybackslash}p{0.62\textwidth}@{}}
\toprule
result & M / T / P\\
\midrule
r51.1(a) & the tachyon map relates M and T and fixes every sector: charges frame-free between M and T; P not evaluated\\
r51.1(b), (c) & statements on the units, the same in M and T\\
\bottomrule
\end{tabular}
\end{center}

\subsection*{v4.9.0-r51 register entries}
\label{subsec:v490-r51-register}

\begin{center}
\small
\begin{tabular}{@{}>{\raggedright\arraybackslash}p{0.34\textwidth}>{\raggedright\arraybackslash}p{0.58\textwidth}@{}}
\toprule
statement & status after v4.9.0-r51\\
\midrule
\(e_7\) as reference axis of the frames & in part: for the tachyon map on all of \(\mathbb O\), for all three maps on \(\operatorname{span}(1,e_7)\) (r51.1)\\
charges in matter and tachyon frame & frame-free; charge conjugation is the time-reversal part \(T\)\\
``frame \(T\) = charge conjugation'' (r41, r43, r47--r49) & means time reversal \(T\) (port exchange), not the tachyon frame; corrected here\\
\bottomrule
\end{tabular}
\end{center}

\subsection*{v4.9.0-r51 ledger and guard update}
\label{subsec:v490-r51-ledger-guard}

\[
R_{\rm fa}^{v490\text{-}r51}=(1,1,0)
\]
for (action of the frame maps on the sectors settled; inner counterparts settled; role of the photon frame \(P\) not evaluated).

The layer spends no guard:
\[
R_{\rm glob}^{(v490\text{-}r51,+)}=(0,1,1,1,1),
\qquad
\Gamma_{\rm guard}^{v490\text{-}r51}=(1,1,1,1,1),
\qquad
\beta_4=1/256\ \text{unspent}.
\]
It does not imply exact finite \(C5\) zero closure, exact all-readout terminality,
\(C6/T2\), a global photon-clock/RH theorem, or global all-readout closure.
